\section{Discussion}
\label{sec:discussion}

Taken together, the results point to a layered interpretation of Brazilian equity dependencies, in line with the multi-scale perspective advocated by Raddant and Di Matteo~\cite{raddant_dimatteo_2023}. The first layer is a market-mode component. It is reflected in the moderate positive average correlation ($\bar{\rho}\approx0.24$), the increase in rolling average correlation around stress periods and the very large leading eigenvalue ($\lambda_1=21.65$, approximately 37.3\% of the trace of the correlation matrix). We interpret this component as broad market synchronization. In the Brazilian context, such synchronization may reflect common exposure to domestic monetary conditions, exchange-rate dynamics, country risk and global commodity cycles, although these channels are not separately identified in the present analysis. The component is also methodologically important because it can dominate raw networks and obscure weaker local structures.

The second layer consists of candidate group and sector structure. Within-sector correlations are higher than between-sector correlations on average, and the RMT decomposition identifies four additional eigenvalues above the random upper edge. Eigenvector loadings and group-mode matrices suggest that commodity, financial, telecommunications and utility-related structures contribute to the non-random part of the spectrum. This reading is descriptive, not causal: we identify structure in the dependency matrix, not the economic mechanisms that generate each edge.

The third layer is noise-compatible residual dependence. Not every correlation is economically meaningful, and not every edge retained by a graph is a stable economic relationship. This is why we combine RMT, heatmaps, dendrograms and network metrics. A result is more credible when it appears across several representations, not only in one visualization.

The network results show that market-mode removal changes structure in a systematic way. Original networks have stronger average edge correlations, while group-mode networks have much lower average edge correlations. However, the group-mode PMFG has higher clustering (0.6653 versus 0.5273). Notably, lower average edge strength does not mean absence of structure. It means that the dominant common component has been removed and localized organization becomes easier to observe \cite{tumminello_2005_tool,tumminello_2010_correlation}. The PMFG is well suited for this analysis because it preserves richer graph structure than the MST, retaining 168 edges with 166 triangles and 55 4-cliques that allow hub-rank shifts, clustering changes and subsector dependency maps to be studied in detail.

Clustering summarizes hierarchical organization, but it does not identify graph-theoretic hubs, bridges or local cycles. The transition from dendrograms to filtered financial networks adds information. Hub reorganization after RMT filtering is meaningful: original hubs such as GGBR4 and BBDC4 are central in networks that include the common component, while group-mode hubs such as GUAR3 occupy central positions after the leading component is removed. This suggests that centrality in financial networks depends on the matrix used to construct the network.

The forecasting results require caution. Machine-learning and simple deep-learning models do not produce uniform gains across horizons and metrics. Instead, network-augmented Random Forest models are competitive at the 5-day horizon, while Ridge regression and CNN-1D are competitive at the 20-day horizon. Feature-importance results indicate that lagged volatility variables remain dominant, with selected PMFG variables contributing smaller but nonzero information. The net interpretation is modest: topology adds structural information from the cross-sectional geometry of the market to the temporal dynamics of individual asset variance.

The broader contribution is methodological. We connect econophysics filtering, financial networks and volatility forecasting in a single reproducible workflow for B3 equities. By documenting each step transparently and evaluating network features through a chronological machine-learning and deep-learning design, the workflow offers a template that can be extended to other emerging markets, different time scales and more sophisticated forecasting architectures.
